\documentclass[acmtods,acmnow]{acmtrans2m}
%&t&{\tt #}&
%&v&\verb|#|&

\acmVolume{32}
\acmNumber{1}
\acmYear{07}
\acmMonth{March}

\usepackage{url}
\newcommand{\tods}{{\em TODS\/}}

\markboth{R.~T.~Snodgrass}
         {Editorial: Single- Versus Double-Blind Reviewing}

\title{Editorial: Single- Versus Double-Blind Reviewing}
\author{Richard T.~Snodgrass\\University of Arizona
}
\begin{abstract}
This editorial analyzes from a variety of perspectives the controversial
issue of single-blind versus double-blind reviewing. In single-blind
reviewing, the reviewer is unknown to the author, but the identity of the
author is known to the reviewer. Double-blind reviewing is more symmetric:
the identity of the author and the reviewer are not revealed to each other.
We first examine the significant scholarly literature regarding blind
reviewing. We then list six benefits claimed for double-blind reviewing and 21 possible
costs. To compare these benefits and costs, we propose a double-blind policy
for \tods\ that attempts to minimize the costs while retaining the core
benefit of fairness that double-blind reviewing provides, and evaluate that policy against each
of the listed benefits and costs. Following that is a general discussion
considering several questions: what this has to do with \tods, does bias
exist in computer science, and what is the
appropriate decision procedure? We explore the ``knobs'' a policy design can
manipulate to fine-tune a double-blind review policy. This editorial ends
with a specific decision.
\end{abstract}

\category{H.2}{Database Management}{General}

% \terms{}

\keywords{anonymous citation, blinding efficacy, double-blind review,
gender bias, single-blind review, status bias}

\begin{document}

\begin{bottomstuff}
Author's address: Richard T.\ Snodgrass, Department of Computer Science,
University of Arizona, Tucson, AZ 85721-0077, {\tt rts@cs.arizona.edu}.

This is a preliminary release of an article accepted by ACM
Transactions on Database Systems. The definitive version is currently
in production at ACM and, when released, will supersede this version.

\end{bottomstuff}

\maketitle

\section{Introduction}

The peer review process is generally acknowledged as central to the
advancement of scholarly knowledge. It is also vital to the
advancement of individual careers.

With so much at stake, it is important to examine, and re-examine, issues
pertaining to review quality on an on-going basis. Thus it is appropriate
that controversy has arisen in our field pertaining to the practice of
double-blind reviewing. ``As scientists, we should rather welcome
all occasions to reflect on the act of writing, evaluating, editing and
publishing research findings. The issue of double-blind refereeing, which
recurs periodically in scientific circles, provides us with such an
opportunity''~\cite[page 324]{genest93}.

Most database journals employ single-blind reviewing, in which the reviewer is
unknown to the author, but the identity of the author is known to the
reviewer. Others employ double-blind reviewing, in which the identity of the
author and the reviewer are not revealed to each other. The arguments for
double-blind reviewing are that it is more fair and that it produces higher
quality reviews.  The arguments advanced against double-blind reviewing
include that it has little effect, that it makes it more difficult for
reviewers to comprehensively judge the paper, and that it is onerous to
administrate~\cite{ceci84}.

The editorial first examines the now substantial scholarly
literature regarding blind reviewing. This literature includes empirical
studies from biomedicine, communication, computer science, economics,
education, medicine, public health, physics, and psychology, retrospective
analyses from computer science, ecology, economics, and medicine, and a
quantitative meta-analysis from psychology. It is useful and instructive to
learn what other disciplines, using diverse approaches, have discovered
about blind reviewing.

Later sections provide a comparison of costs and benefits, a proposal for
a double-blind reviewing procedure, and an general discussion of relevant
questions. We end with a specific decision.

\section{An Analysis of the Literature}
In the following\footnote{This section contains a subset of the literature analysis that
appeared in a companion paper~\cite{DBR}, with an added discussion of
psychological studies of bias.}, we first define the various terms used in the
literature. We then examine in some depth
the general issues of fairness. We end with a brief summary of this complex
sociological question.

\subsection{Terminology}

ACM defines a {\em refereed journal} or {\em refereed conference} as one
that ``is subjected to a detailed peer review, following a defined, formal
process according to a uniform set of criteria and
standards.''~\footnote{ACM Policy on Pre-Publication Evaluation, at
\url{http://www.acm.org/pubs/prepub_eval.html}} Material appearing in such
venues is distinguished from
{\em formally reviewed material} (``subjected to a structured evaluation and
critique procedure following a defined process uniformly applied as with
refereeing, only without requiring that the tests of scholarly originality,
novelty and importance be applied''), {\em reviewed} (``subjected to a more
informal and not necessarily uniform process of volunteer review, with
standards dependent upon the publication and the type of material''), {\em
highly edited} (``professionally edited, usually by paid staff, with primary
emphasis on exposition, graphic presentation, and editorial style rather
than on content and substance''), and {\em unreviewed} (``published as
submitted, with or without copyediting''). ``Reviewing'' in the present
document refers to peer review for a refereed journal or conference.

{\em Peer review} is the use of predetermined reviewers, in the case of program
committees, or ad hoc reviewers, in the case of reviewers
for most journals, who individually read the submitted
manuscript and prepare a written review. Sometimes, as in the case of some
conference program committees, reviewers will
subsequently either physically or electronically meet to discuss the papers
to arrive at an editorial decision. For most journals, the Associate Editor
handling the paper or the Editor-in-Chief will make the final editorial
decision.

In the vast majority of refereed database conferences and journals, the
identity of the reviewer(s) is not revealed to the author(s), ostensibly to
ensure more objective reviewing. This is termed {\em single-blind reviewing}
or, less frequently, {\em one-eyed review}~\cite{rosenblatt80}. (The terms
{\em reviewer} and {\em referee} are used interchangeably in the literature.)

In an effort to achieve more objective reviewing, a venue can also request
that the identity of the author be removed from the submitted manuscript, a
process termed {\em blinding the manuscript}. When the identity of the
authors and their institutions is kept from the reviewers, this is termed
{double-blind reviewing}. Note that the Editor-in-Chief and Program Chair,
and generally the Associate Editor, are made aware of this information via a
separate cover sheet not shared with the reviewers.

The psychological sciences utilize a different terminology that conveys a
subtle philosophical shift. When the identities of the authors and reviewers
are not revealed to each other, it is termed in these sciences a {\em
masked} reviewing process. Note the symmetry of this terminology. The
American Psychological Association {\em Guide to Preparing Manuscripts for
Journal Publication}~\cite{APAGuide} states, ``Peer review is the backbone
of the review process. Most APA journals, like the majority of other
professional publications, practice anonymous, or masked, reviews. Authors
and reviewers are unaware of each other's identities in most instances, an
arrangement designed to make the process more impartial.'' The implication
is that revealing either the reviewer's identity or the author's identity
breaks the mask. Presumably single-blind reviewing would then be termed
``non-masked,'' but the APA doesn't use the term. (The term
{\em unmasking} denotes revealing the identity of a reviewer to a
co-reviewer~\cite{vanRooyen99}; we don't consider that practice here.)

This editorial will use the terms single-blind and double-blind
reviewing, as well as their respective three-letter acronyms, {\em SBR} and
{\em DBR}.

Venues differ in who does the blinding/masking of a submission. We will use
the term {\em author masking} when the author removes identification from the
paper before submitting it and {\em editorial masking} when such
identification (generally, author name and affiliation) is removed in the
editorial process before sending the manuscript to the reviewers. Procedures
differ in how aggressive is the required author masking and the actual
editorial masking. Self-citations and other first-person references in the
body of papers are generally retained in editorial masking.  While author
masking can be more thorough, because authors would know what kind of
information is revealing, authors through various devious means can
circumvent both kinds of masking.

Previous literature reviews (see the companion paper~\cite{DBR} for a survey)
emphasized three primary aspects relevant to blind reviewing: fairness to
authors (to unknown authors or to authors affiliated with unknown
institutions, to less-published or to proficient authors, to both genders),
review quality, and blinding efficacy. The following section addresses
fairness; the companion paper contains a comprehensive survey of the other
two aspects.

\subsection{Fairness}\label{sec:fairness}
The fundamental argument for double-blind reviewing is that it is more fair to
authors (and thus, indirectly, to readers). The argument proceeds as
follows: The judgment of whether a paper should be accepted for publication
should be made on the basis of the paper alone: is what the submission
states correct, insightful, and an advancement of the state-of-the-art? The
editorial judgment should not be made on extenuating circumstances such as
who wrote the paper or the professional affiliations of the authors. By
blinding the submission, the reviewers cannot take these peripheral aspects,
which are not relevant, into account in their review.

Psychology researchers have studied bias in judgment.
Expectations are one source of bias. Quality is judged relative to
expectations, which is why a \$6 cup of coffee tastes better than a \$2 cup
of coffee, holding everything else constant. This sort of bias is both
conscious and unconscious. People know that they expect more expensive
coffee to taste better, but they don't realize that the mere fact of having
that expectation actually makes it taste better~\cite{shiv05}. Some
reviewers have an expectation that ``papers from $\ldots$ are never any
good.''

There is now also a good amount of evidence from psychological
studies that judgment often operates in an unconscious
fashion~\cite{greenwald95}. People who don't {\em know} that they are biased
tend to be those who are the most biased~\cite{wilson94}. And paradoxically,
under some conditions, experience and motivation, both of which are in ample
supply in reviewing, can accentuate some forms of judgment
bias~\cite{kardes05}: ``people believe that they learn a lot from experience
even when experience is actually irrelevant'' [ibid, pp.~149--150]. Such
experience could involve the identity or affiliation of the author.

The relevant question then is whether these conscious and unconscious sources
of bias actually come into play in reviewing of conference or journal articles.
The analysis of fairness in the extant literature concerns (a) fairness to
unknown authors or institutions, (b) fairness to prolific or to less-published authors, and
(c) gender equity. There is also the related issue of the perception of
fairness. The following sections will elaborate on each of these concerns.

\subsubsection{Fairness to Unknown Authors or Institutions}

Some evidence from retrospective and experimental studies suggest that when
the authors' names and affiliations are known, reviewers may be biased
against papers from unknown authors or institutions, termed {\em \hbox{status}
bias}~\cite{cox93}. We now examine the studies that attempt to detect
status bias, in chronological order.

%replace with quotes from actual article (Crane)
A retrospective study of manuscripts that had been submitted to {\em The Physical Review}
between 1948 and 1956 found that ``some 91 per cent.~of the papers by
physicists in the foremost departments were accepted as against 72 per
cent.~from other universities''~\cite[page 85]{zuckerman71}. Two possible
explanations were offered: status bias and ``differences in the scientific
quality of the manuscripts coming from different sources''~[ibid].  Crane
[1967] provided an alternative explanation, ``that common viewpoints rather
than personal ties could explain acceptance''~\cite[page~227]{dalton95}, yet
such common viewpoints would have been present in both SBR and DBR.

An early experiment found that ``the effect of institutional prestige failed
to attain significance in any one of the measures''~\cite[page
70]{mahoney78}.  ``Experimental manuscripts were sent to 68 volunteer
reviewers from two behavioristic journals. ... Institutional affiliation was
also manipulated on the experimental manuscripts, with half allegedly
emanating from a prestigious university or a relatively unknown
college''~[ibid].

Another retrospective study, this of the records of reviews of a society
which publishes research journals in two areas of the physical sciences,
found large differences in how papers from minor and major universities are
reviewed: ``minor university authors are more frequently evaluated
favourably (ie less critically) by minor university referees, while major
university authors are more often evaluated favourably by major university
referees than they are by those affiliated to minor universities. It would
therefore appear that when referees and authors in these areas of the
physical sciences share membership of national or institutional groups, the
chances that the referees will be less critical are increased. ... Personal
ties and extra-scientific preferences and prejudices might, of course, be
playing a part as well. But it appears that, even in the absence of these
personal factors, the scientific predispositions of referees still bias them
towards less critical evaluation of colleagues who come from similar
institutional or national groups, and so share to a greater extent sets of
beliefs on what constitutes good research''~\cite[pp.~274--5]{gordon80}.

A seminal experiment~\cite{blank91} demonstrated status bias in reviewing
more directly.  In this experiment, every other paper that arrived at the
{\em American Economic Review} was designated as double-blind. For these
papers, an editorial assistant removed the name and affiliation of the
author from the title page and typically scanned the first page for
additional titles or notes that would identify the author (that is, those
manuscripts were editorially masked). This experiment lasted for two years.

The relevant issue was ``whether the ratio of acceptance rates between
institutional ranks in the blind sample differs from the corresponding ratio
in the nonblind sample.''~\cite[page~1053--1054]{blank91}. It was found that
this ratio did not differ for those at top-ranked departments and those at
colleges and low-ranked universities. All other groups, in that important
gray area where editorial judgment is most needed, had substantially lower
acceptance rates in the blind sample than in the nonblind sample; in some
cases, the acceptance rate dropped by more than 7 percentage points.  She
found similar differences with referee ratings between SBR and DBR.

A retrospective study of single-blind
reviews for the {\em Journal of Pediatrics} and published in {\em JAMA}
found only partial evidence for status bias, that ``for the 147 brief reports, lower institutional rank was
associated with lower rates of recommendation for acceptance by reviewers
($P<.001$). ... For the 258 major papers, however, there was no significant
relationship between institutional rank and either the reviewer's
recommendations ($P$=.409) or the acceptance rate ($P$=.508)''~\cite[page
138]{garfunkel94}.

Another retrospective analysis of single-blind reviews also published in
{\em JAMA} found evidence of
status bias at a coarse geographical level~\cite{link98}. In this
analysis of original research articles submitted to {\em Gastroenterology}
during 1995 and 1996, it was found that ``reviewers from the United States
and outside the United States evaluate non-US papers similarly and evaluate
papers submitted by US authors more favorably, with US reviewers having a
significant preference for US papers''~\cite[page~246]{link98}.

The experimental evidence is mixed concerning status bias present for
top-ranked authors and institutions. The evidence is quite compelling that
status bias is possible, perhaps prevalent, in SBR for most other authors
and institutions, presumably for those papers most needing the critical
evaluation of reviewers.

\subsubsection{Fairness to Prolific Authors}

There have been several studies that have looked at the impact of blinding
on prolific authors; these studies were discussed in the companion
paper~\cite{DBR}. Contradictory results from these studies render it
impossible to say anything definitive about the impact of blinding on
prolific authors. However, there does seem to be evidence of some kinds of
bias with SBR.

\subsubsection{Gender Equity}

When reviewers know the identity of the author(s) of the submitted
manuscript, gender bias is also a possibility. Several disciplines have
launched in-depth studies based on concerns of gender equity.

%- Blackburn\&Hackel refs Fella Salla\&Gragman02,
%Kelicie\&Merckelbach02, and Wenneria\&Wold97
% Billard cites [Billard 91-94], check out her papers

Blank's experiment, described earlier, was in fact initiated due to concerns
of gender bias. The {\em American Economic Review} journal had employed SBR
for most of its recent history, except during a period of 1973--1979 when
the then-current editor adopted DBR~\cite{borts74}. In the mid-1980's, the
American Economic Association's Committee on the Status of Women in the
Economics Profession formally expressed its concern about ``the potential
negative effect on women's acceptance rates of a single-blind
system''~\cite[page~1045]{blank91}. As a result, Blank was asked by the
current editor of the {\em AER} and the Board of Editors to design and run a
randomized experiment looking into this potential effect.

Due to the careful randomization design of this experiment, one can compare
acceptance rates between the blind and nonblind samples, and indeed, there
were striking differences. ``For women,
there is no significant difference in acceptance rates between the two
samples. For men, acceptance rates are significantly
higher in the nonblind sample.''~[ibid, page~1053]. When reviewers knew that
that paper was authored by a male, they accepted a higher percentage (15\%,
versus 11\%) than if the paper was blinded. ``One can compare
acceptance rates between the blind and nonblind samples without other
control variables because the randomization process guarantees that papers
by women (and men) in each sample have identical distributions of
characteristics''~[ibid].

Blank emphasized the core issue: ``whether the {\em ratio} of male to female
acceptance rates in the nonblind sample is different from that in the blind
sample. In both samples, women's acceptance rates are lower than men's, but
the differential in the blind sample is smaller. While women in the blind
sample have an acceptance rate only 1 percentage point below that of men,
their rate is 3.8 percentage points lower in the nonblind sample''
[ibid]. Here the results were statistically insignificant, perhaps because
there were too few observations of papers authored by women.

Would DBR result in a large increase in acceptances of papers by women?
``While there is some indication in these data that women do slightly better
under a double-blind system, both in terms of acceptance rates and referee
ratings, these effects are relatively small and statistically
insignificant. Thus, this paper provides little evidence that moving to a
double-blind reviewing system will substantially increase the acceptance
rate for papers by female economists''~[ibid, page~1063]. Interestingly, the
{\em American Economic Review} reverted to DBR after Blank' study was published.

The Modern Language Association's (MLA) experience was striking:
going to DBR resulted in a large increase in acceptances by female authors.
``Contributed papers at MLA meetings had first to survive a
review stage before acceptance to be read. Prior to 1974, these papers were
refereed with the author's name intact. In 1974, double-blind refereeing was
tried with the effect that the number of women and of new investigators
having papers accepted doubled from previous years. This number doubled
again when repeated in 1975, until, by 1978, the proportion of acceptances
among women and new researchers was comparable to that for men. The MLA Board
subsequently decided in 1979 to use double-blind refereeing for all their
publications''~\cite[page~321]{billard93}. The impetus for this change was
the perception of gender bias. ``A number of women complained to the Modern
Language Association in the United States that there were surprisingly few
articles by women in the association's journal, compared to what would be
expected from the number of women members. It was suggested that the review
processes were biased. The association vigorously denied this but under
pressure instituted a blind reviewing procedure under which the names of the
authors and their institutional affiliations were omitted from the material
sent to the reviewer. The result was unequivocal: There was a dramatic rise
in the acceptance of papers by female authors''~\cite[page 217]{horrobin82}.

It is possible that the small observed effect in Blank's study (in contrast
to the MLA experience) was due to the low number of submissions by women to
{\em AER}. Certainly the computer science field is closer to economics than
modern languages in its participation of women.

These studies show that revealing author identity, specifically the gender of
the author, can sometimes have an effect on acceptance rates.

\subsubsection{The Perception of Fairness}

A {\em perception} of possible bias may be just as damaging as actual bias.

The Institute of Mathematical Statistics (IMS) New Researchers' Committee
(NRC) report stated, ``The NRC feels that the current system [SBR] has the
potential for bias or perceived bias against NRs [new researchers], women
and identifiable minorities, (a disproportionate number of the latter two
categories are NRs)''~\cite[page~165]{altman91}.  In a response to
discussants of that report, the NRC reasserted a year later, that ``much of
the value of double-blind refereeing lies in the community perception of
fairness''~\cite[page~266]{altman92}.

The experience with this controversy at the IMS indicated a split between
new researchers, which ``strongly endorses double-blind refereeing. ... It
seems likely that [this] represents the majority opinion among new
researchers, although support for double-blind refereeing is not unanimous
among new researchers'' and senior members: `` `negative but sympathetic'
... seems to be a majority view among those senior enough to have been
involved in the editing process''~\cite[page~311]{cox93}. However, a survey
to IMS members ``indicates strong support for double-blind refereeing in the
IMS journals''~[ibid].

A responder to the IMS report~\cite{cox93} stated, ``Refereeing is
{\em perceived} by many writers as being subject to various kinds of biases:
biases in favor of male or female, young or established, national or foreign
researchers, working at small or large institutions, in well-developed or
developing countries and so on. Whether such biases are sufficiently strong
and widespread to distort the whole review process is beyond the point. So
long as the {\em potential\/} for abuse is there, we should guard against it,
and double-blind refereeing is but one means of ensuring such
protection''~\cite[page 324]{genest93} (emphasis in original).

\subsection{Quality of Reviews}

From the analysis in the companion paper~\cite{DBR}, one must conclude that
the jury is still out. It has not been shown convincingly that either SBR or
DBR can, by revealing or by hiding the identity of the author and
institution, increase the quality of the reviews of a submitted manuscript.

\subsection{Efficacy of Blinding}\label{sec:efficacy}

The companion paper~\cite{DBR} examined a number of studies of blinding
efficacy. Blank's conclusions apply to the many studies generally. ``On
the one hand, a substantial fraction---almost half---of the blind papers in
this experiment could be identified by the referee. This indicates the
extent to which no reviewing system can ever be fully anonymous. On the
other hand, more than half of the papers in the blind sample {\em were}
completely anonymous. A substantial fraction of submitted papers are not
readily identified by reviewers in the field. ... Those blind papers that
are correctly identified by the referees ... are skewed in favor of authors
who are better known or who belong to networks that distribute their working
papers more widely''~\cite[pp.~1051--2]{blank91}. Ceci and Peters
conclude that ``Although there are occasional lapses in the preparation of
manuscripts by authors and failures to screen manuscripts by editorial
staff, we are impressed by the overall efficiency of blind
review''~\cite[page 1494]{ceci84}.

\subsection{Summary}
``There is a long tradition attached to the peer review system. As {\em
users} of science, we all depend on it: our professional realizations are
based upon the work of others, and we count on journal (and book) editors to
separate the wheat from the tares. Although there is no such thing as
perfection, it would be a disservice to the profession if too many
scientific writings addressed irrelevant issues or contained gross factual
errors. As {\em producers} of science, it is also in our interest that the
system be fair: favoritism, discrimination and condescension bring
discredit on the entire operation and ultimately work against the
discipline, even if individual benefits occasionally may accrue in the short
term''~\cite[page 324]{genest93}.

We have attempted here to summarize the many studies of the varied aspects
of blind reviewing within a large number of disciplines.

Concerning the central issue of fairness, Blank's summary in 1991 of the
literature still holds true fifteen years later.  ``In summary, the
literature on single-blind versus double-blind reviewing spans a wide
variety of disciplines and provides rather mixed results. Few of the
empirical tabulations provide convincing evidence on the effects or
non-effects of refereeing practices, largely because of their inability to
control for other factors in the data. If not fully convincing, however,
there is at least a disturbing amount of evidence in these studies that is
consistent with the hypothesis of referee bias in single-blind
reviewing''~[page~1045]. Many studies provide evidence that DBR is more fair to
authors from less-prestigious institutions and to women authors.
%replace with direct quote
%One paper~\cite{mcgiffert88} concludes ``Double blinding is generally
%considered to benefit the scholarly dis\-advantaged--i.e., women, minorities,
%and authors at universities that lack prestige''~\cite[page~236]{dalton95}.
Such differences are likely to matter even more for highly-selective
conferences and journals.

A companion paper~\cite{DBR} looked at related issues.  Concerning quality
of reviews, it is not known definitely whether either SBR or DBR results in
a higher quality of reviews.  Most of the studies discussed here and in the
companion paper utilize editorial blinding, which has been shown to be
successful about 60\% of the time, across many disciplines. Removing text
allusions and self-citations would increase success rates to perhaps 75\%.
The prevalence of DBR has increased dramatically over the last fifteen
years, to the point where most scientific journals now employ double-blind
reviewing.

\section{An Analysis of Costs Versus Benefits}

As the previous literature survey emphasized, there are demonstrable, albeit
intangible {\em benefits} of a double-blind reviewing process. And it
emphasized as well that there are very real and tangible {\em costs}, to the
journal as well as to the author and perhaps to the reviewer, of adopting
DBR. These costs depend in part on how the blinding is done, with the
efficacy directly related to the effort expended. These costs also vary
depending on the needs of the scholarly discipline covered by the journal in
question and on the culture of the scholarly community served by that
journal.

Journals strongly desire to
fairly evaluate submitted manuscripts, while simultaneously keeping costs in
control. The policy question before each journal and each scholarly
publisher is thus the following. Is the documented benefit of equity worth
the administrative cost? At what price fairness?

We now enumerate the benefits and costs of adopting DBR and
proposes a specific procedure for \tods\ that attempts to
minimize those costs while retaining the documented benefits.

\subsection{Perceived Benefits of DBR}\label{sec:benefits}

We first list benefits claimed for DBR.

\begin{description}
\item [DBR avoids a bias towards top-ranked authors and institutions]
Perhaps SBR encourages an unconscious bias towards prominent authors or institutions.

\item[DBR avoids a bias towards prolific authors] Prolific authors will be
well-known, and may be treated more gently by reviewers.

\item[DBR encourages higher-quality reviews] Perhaps by not knowing the
identity of authors or their institutions, reviewers will write more
authoritative reviews.

\item[DBR is more fair to women authors] Studies have shown that revealing the
gender of the author can have an effect on acceptance rates.

\item[DBR is more fair to non-top-ranked authors and institutions] The
evidence is quite compelling that status bias is possible, perhaps
prevalent, in SBR for such authors.

\item[There is a perception by some that DBR is more fair] 
Several task forces and committees have called for DBR.
\end{description}

\subsection{Perceived Costs of DBR}\label{sec:costs}
There are very real costs to DBR that must be objectively
considered when contemplating a change in a journal's editorial masking
policy. Some of these costs apply to most scholarly journals; other costs
are specific to a computer science journal.

\begin{description}
\item[Quality of reviews is reduced] Perhaps DBR, by
depriving the reviewer of the knowledge of the identity of the author(s),
their institution, their prior work, or the biases identifiable through
their prior work, reduces the quality of reviews.

\item[Intentional bias is reduced] Perhaps DBR disallows intentional bias
that attempts to counter unconscious bias by other reviewers.

\item[It is difficult to mask submitted manuscripts]
Editorial masking requires either blackening out identifying information
from a paper (which is difficult if the paper is submitted electronically)
or requires that the source be also submitted so that it can be
edited. Author masking requires effort on the part of the author to remove
identifying information.

The web can be searched by reviewers to find related
papers (especially those that are similarly titled by the same author, such
as technical reports and related papers), thus making it fairly easy to
narrow down the set of suspected authors.

\item[DBR increases the probability of non-novel papers getting accepted]
Perhaps DBR will discourage reviewers from looking up related work, and so
may artificially inflate the perceived contribution of the work.

\item[It is difficult to fully mask systems papers]
A mode of inquiry common in computer science is the creation of often
extensive software artifacts, often over a period of several
years. In this research methodology, a succession of papers is produced,
each of which builds on the artifacts and prior papers. Thus it is difficult
to impossible to fully mask manuscripts coming out of such projects, as the
contributions are necessarily incremental.

\item[DBR jeopardizes ``power projects'']
There is a history of high-impact projects in computer science; examples
include the Berkeley RAID project, the Cornell Program Synthesizer, Berkeley
Ingres,
IBM's System R, and UCSD Pascal. Such projects may not be compatible with
DBR. Papers from these projects benefit from the ``brand'' that has been
created by prior distribution of papers and software artifacts such as freely
distributed source code; DBR does not allow this brand to be exploited or
even acknowledged by the submitted manuscript.

\item[Forward-looking systems work, in particular, suffers from
anonymization]
Such systems work can get obfuscated if the author tries to truly anonymize
it, or it is really not anonymized.

\item[Anonymization is a distraction to students learning to write  
clearly]
Do we want to emphasize to students how to make something (in this case,
authorship) unclear?

\item[Some journal policies preclude DBR]
As one example, \tods\ requires that submissions invited from conferences
utilize some of the original reviewers from that conference's program
committee in the review of the journal submission. DBR would
preclude inviting papers in the future from conferences.

\item[DBR precludes expanding a conference paper into a journal submission]
A submission that is an expanded version of a conference paper, containing
much of the same prose and often a similar if not identical title as the
conference paper, cannot be easily masked.

\item[DBR makes it difficult for conflicts of interest to be detected] How
can a reviewer determine whether a conflict of 
interest exists with an author of a blinded manuscript?

\item[DBR prevents the dissemination of research results before submission]
Strictures on prior submission in DBR policies could prevent working papers and
technical reports from being distributed for comment, because such
dissemination could reduce the efficacy of author masking.

\item[DBR gets papers rejected for failing to sufficiently cite the author's
own work]\ 
Citations removed during the blinding of a submitted manuscript may cause a
paper to be rejected because it does not cite relevant work.

\item[DBR increases the possibility of plagiarism] With DBR, a reviewer
might assume material is from one author and judge it legitimate, even if
the material was actually plagiarized from another author.

\item[DBR raises the possibility of self-plagiarism] With DBR, other work
by the authors may be anonymously cited or not cited at all, making it more
difficult for reviewers to catch self-plagiarism.

\item[DBR places extra burdens on reviewers] DBR, through the
anonymization process, may render the paper harder to understand or to review.

\item[DBR decreases enthusiasm of reviewers] Perhaps DBR (and even the
discussion about DBR) sends the message that the journal does not trust its
reviewers.

\item[DBR reduces timeliness]
If the enthusiasm of reviewers is reduced, or if the process is cumbersome,
timeliness may be adversely affected.

\item[DBR places extra burdens on the Editor-in-Chief and Associate
Editors]
Depending on how the mechanism for ensuring DBR is designed, this mechanism
may impose more work on the Editor-in-Chief (EiC) or on the Associate
Editors (AEs).

\item[DBR places extra burdens on the authors]
Similarly, depending on how the mechanism for ensuring DBR is designed, this
mechanism may impose more work on the authors.

\item[SBR versus DBR should not be a local decision]
There is the ``cost'' of incompatibility between the editorial policies of
various journals serving a scholarly community. It is less confusing if these
journals present a consistent set of policies. Otherwise a paper will have
to be altered if it is rejected from one journal and subsequently submitted
to another journal.
\end{description}

\subsection{A Proposal for a Double-Blind Policy}\label{sec:proposal}
To compare the costs and benefits of DBR, we propose a double-blind policy
for \tods\ that attempts to minimize the costs while retaining the benefit
of fairness.\footnote{This proposal was hammered out with the \tods\
  Editorial Board over an intensive two-month discussion via email.}

This proposal roughly follows the approach utilized by the ACM {\em SIGMOD}
conference, which has been using DBR since 2001. That conference's
double-blind policy specifics and procedures have been refined over the past
six years, and have been applied to several thousand submissions; this
policy is now well-known in the database community.  An
informal check of several other conferences and journals utilizing DBR
indicated that the instructions from the {\em SIGMOD} conference are by far
the most specific and helpful in their guidance to authors.

For each aspect (applicability, submission, reviewing, revision, acceptance,
transition), we state the proposed policy, then review the rationale behind
that proposal. The following section revisits the benefits and costs of
adopting this policy.

As mentioned, two kinds of masking have been discussed in the prior literature: editorial
blinding and author blinding. In general, there is a wide spectrum between
the extremes of a manuscript that explicitly states its authorship (unmasked)
and one for which there is no possibility of a reviewer discovering the
identity of any of the authors (complete masking). Indeed, the latter
situation is not feasible. There is always the potential of a reviewer
discovering, through a web search, through conversation with colleagues,
through knowledge of the literature, or through other means, the identity of
one of the authors. But between the points of {\em maximal realizable
masked} and unmasked is a plethora of levels.

Studies have shown that status and gender bias are reduced even when author anonymity is
rather compromised. Most studies utilized editorial masking, which is
successful a little more than half the time. This suggests that one can move
along this spectrum, increasing (or reducing) the benefits of DBR while also
increasing (or reducing)
the cost. The challenge is to figure out where the aggregate of {\em benefit
minus cost} is maximized.

The following proposal attempts to get that balance right. It acknowledges
that submissions with less masking have a greater chance of having one or
more authors revealed, which might introduce or increase status or gender bias during
the review. But also acknowledged is that {\em uncertainty} as to the
authors' identity is often sufficient to realize most or all of the benefits of
masking.

\subsubsection{Applicability}\label{sec:applicability}
Reviewing of all submissions to \tods\ will be double-blind.
All submissions to \tods\ must be masked by the author(s), following
the simple instructions given in the proposed author guidelines
(Appendix~\ref{app:authors}).

\noindent
{\sffamily Rationale}

Author masking avoids the procedural difficulties of editorial masking,
distributes the workload, and, most importantly, significantly increases the
efficacy of masking. The literature analysis in the companion paper found
that ``most of the studies discussed here utilize editorial blinding, which
has been shown to be successful about 60\% of the time, across many
disciplines. Removing text allusions and self-citations would increase
success rates to perhaps 75\%.''~\cite[page 18]{DBR} We note that the {\em
SIGMOD} conference also utilizes author masking.

Editorial blinding has been shown to be effective in the past, and more
aggressive author blinding is even more so. Blank concluded ``On the one
hand, a substantial fraction---almost half---of the [editorial] blind papers in this
experiment could be identified by the referee. This indicates the extent to
which no reviewing system can ever be fully anonymous. On the other hand,
more than half of the papers in the blind sample {\em were} completely
anonymous. A substantial fraction of submitted papers are not readily
identified by reviewers in the field.''~\cite[pp.~1051--2]{blank91} No
masking process can be totally successful, yet many of the benefits of DBR
accrue from uncertainty in the identity of the author.

This analysis is based on scientific studies that were performed before the
advent of the Internet and search engines. In the studies, about half the reviewers
were able to guess the identity of at least one of the reviewers. This is
partially because the studies utilized editorial masking. So while author
masking does better, the availability of web searches decreases the
efficacy, so they balance out somewhat, though there is no good data on the
efficacy nor on the impact of status or gender bias in today's more
highly-connected environment.

The proposed author guidelines are quite specific on how to ensure
anonymity, how and when to use anonymous citations, specifics on papers
involving well-known or unique systems, and anonymity in
revisions.

The rules for anonymous citations requires such citations both for the
authors' own work and that of others, for certain categories of papers
(primarily, those that only the author would be aware of). This was done
originally in the {\em SIGMOD} conference author instructions to ensure that
reviewers could not infer, from what was anonymous and what was not, the
identity of the
authors. Say that the paper referenced two papers that it relied
heavily on or built upon. One is by author $A$ and the other is an anonymous
citation. If that anonymous citation was not masked, but rather was of
author $B$, of a paper to appear, the reviewer could guess that an
author of the submitted paper was $B$, because who else could have built on
work that hasn't even appeared? That is the rationale behind requiring
anonymous citations even of work by others that is to appear.

As noted before, \tods\ policy requires that submissions
invited from conferences utilize some of the same reviewers as the
conference submission. Those reviewers will of course know the identity of
the authors of
the conference paper. However, the other reviewers of the submitted paper
need not be told exactly who the authors are. Anonymizing the rest of the
journal submission also introduces uncertainty in the reviewer's mind. Were
authors dropped or added to this submission? Was the extension carried out
by the senior author or by a junior graduate student?

An alternative to the mandatory policy for the remaining submissions is {\em
elective DBR}, where the author decides whether to blind their
submission. The experience with elective DBR in other journals is that it
works well when there is an established tradition in the community for DBR
(as in American Psychology Association, about half of whose journals employ
mandatory DBR and most of the rest employ elective DBR) and that it is not
invoked often in communities without such a tradition, such as
physics,\footnote{\url{http://forms.aps.org/author/ndbm.pdf}\quad The {\em
IEEE Transactions on Knowledge and Data Engineering} is a more relevant
example. This journal employs elective DBR, which is minimally
used---perhaps 5\%~\cite{werner06}.)} perhaps because authors feel that such
author-requested blinded submissions will be judged more harshly.

That said, there are some papers in computer science and particularly in
database systems that are effectively impossible to blind, and it would be
unfortunate if such papers were disallowed. Note however that the
instructions to the authors do not require complete blinding, which is in
any case not possible for any paper. Rather, all that is required is to
follow the simple six steps in those instructions. For some papers, these steps
will only partially obscure the authorship, which is understood. Uncertainty
of authorship is still preferable to complete knowledge.

\pagebreak
A previous version of this proposal placed a higher bar on the degree of
obfuscation of the submitted manuscript, but also stated three exception
categories:
\begin{itemize}
\item Submissions invited from conferences,
\item Submissions that are extensions of papers that have
been previously published in conferences, and
\item Submissions for which blinding is effectively impossible.
\end{itemize}
As several noted, such categories have the potential to create a
double standard and cause inconsistencies. Also for such papers it is still
important to emphasize to the reviewer that the identity of the author
should not influence their review. Hence, in this proposal the
exceptions were dropped and the standards for masking reduced.

\subsubsection{Submitting a Manuscript}
\tods, as with most ACM journals and transactions, uses ScholarOne's
web-based ManuscriptCentral manuscript tracking system to automate
submission and reviewing.

The manuscript will be submitted normally. The EiC will
examine the manuscript to ensure that it has been properly blinded.

The cover letter and other correspondence from the author(s) will be marked
in ManuscriptCentral (MC) by the EiC as not to be revealed to the
reviewers. The handling AE will have full access to
these materials.

As mentioned in the author guidelines (Appendix~\ref{app:anoncit}), the
cover letter should contain the full details of each anonymous citation as
well as a list of people who constitute a conflict of interest.

\noindent
{\sffamily Rationale}

It is important that the EiC and AE have available the details of the
anonymous citations for the decisions described below.

A list of conflicts of interest in the cover letter will help avoid cases
where a reviewer gets well into a review before realizing that (s)he knows
an author and has a conflict with that author. The names should be
categorized; there have been cases where authors put under CoI reviewers 
perceived to have a high reviewing standard.

The {\em SIGMOD} conference has required, for several years, that conflicts
of interest be stated, in a restricted fashion: one must
indicate when submitting a paper which program committee member(s) have a
conflict of interest, so that the paper will not be assigned to them. The
{\em SIGPLAN} conferences have a similar requirement.

The definition of conflict of interest proposed in the author guidelines is identical to that from
NSF~\cite{NSF}, which is well-thought-out.\footnote{Jennifer Widom
introduced a specific definition of conflict of interest for the {\em
SIGMOD'05} conference; that definition has continued in subsequent
conferences. The {\em SIGMOD} definition differs from that proposed here in that it considers
collaborators for only 24 months to have a conflict, whereas the NSF
guidelines state 48 months. The other difference is that the {\em SIGMOD}
conference definition also includes the condition ``The PC member has been a
co-worker in the same department or lab within the past two years.'' This
was not included in this \tods\ proposal to avoid having the author list all
members in the department or lab. In most cases the AE will be able to
determine this directly from the affiliation of the possible reviewer.}

Having the author prepare a list of conflicts of interest may take an hour
or so. (For those authors who have submitted an NSF grant, satisfying this
requirement will be easy.) Once such a list is created, it is not hard
to maintain.

Some of the advantages of this list apply to SBR as well. Some other
journals have a similar CoI identification requirement; there doesn't seem
to be a strong correlation with single/double-blind reviewing and CoI
identification. 

\subsubsection{Reviewing}
The identity of the
authors will be visible to the EiC and Associate Editor; this meta-data will
not be revealed by MC to the reviewers. Templates of
email to be sent to reviewers will be modified to eliminate identifying
information.

Reviewers should direct questions about a paper referenced by an anonymous
citation within the manuscript (say, arising from a consideration of the
\tods\ novelty requirement) to the AE, who would then
provide the cited paper (or equivalently, a full citation) to the
reviewer.
Instructions to reviewers are given in Appendix~\ref{app:reviewers}.

In cases where an AE has an intransigent referee who does not respond, the AE
sometimes just does the review him/herself, as permitted
in the AE manual~\cite{AEmanual}. While this is strictly speaking a
violation of DBR, the AE will continue to have this option available.

\noindent
{\sffamily Rationale}

By removing all author names and affiliations from email templates, 
and by simply removing author names and affiliations from the web
pages seen by reviewers, these email templates and web pages can be used for
double-blind reviewing.

Anonymous citations are handled similarly to the way they were handled in {\em
SIGMOD'06}, which worked well for the few requests that were made. As each
AE handles only a few \tods\ submissions each year,
responding to reviewer requests should not constitute a burden.

The instructions to the reviewers does not preclude them from doing web
searches or other searches and does not disqualify them if they do discover
one of the author's identity. Attempts to prevent such discovery in every
single case are bound to fail. Rather, reviewers are simply relied upon not
to go to unusual lengths to try to discover the identity of the
author. Again, as the culture of the database community becomes more
familiar and comfortable with the concept and process of DBR, things should
settle out.

Note that the cases where reviews are provided a full citation does not
necessarily reveal the full authorship of the submitted paper. The omitted
author list on the title page still signals to the reviewer not to utilize
that information in their assessment.

\subsubsection{Assessing the Disclosure and Novelty Requirements}\label{sec:assessing}

\tods\ has specific
guidelines\footnote{\url{http://www.acm.org/tods/Authors.html#PriorPublicationPolicy}}
for disclosing related work by the author(s)
of the submitted paper and for what represents adequate contribution over
existing published work.

The above procedures allow the disclosure and novelty requirements to be
checked. It is completely up to the AE to decide which reviewers should be
told about an anonymous citation.

\pagebreak

\noindent
{\sffamily Rationale}

\tods\ policy already requires a {\em prior publication policy prescreen}
by the AE.
The cover letter and the submission itself will provide
sufficient information for the AE to perform the prescreen.

This is especially important for papers that are extended from conference
papers. A recent study, which looked at the papers themselves, found 14
(non-invited) papers extended from conference papers published in \tods\
over calendar 2003--2005~\cite{Montesi06}. There were also 22 papers invited
from conferences (mainly {\em SIGMOD} and {\em PODS}) that were published
during that three-year period. This works out to 56\% of published
papers. However, as invited submissions have a much higher success rate, it
is estimated that slightly less than half of the submissions to \tods\ are
extended from conference papers.

The disclosure requirements (30\%) for papers that are extensions of
conference papers have in the past been generally checked by reviewers. That
will continue under DBR, in two phases. First, the AE will check the paper
to ensure that statements about contributions beyond prior published work
(e.g., ``We added a section on algorithms'', ``we now present a detailed
proof of Theorem 4'') are accurate, consulting with the referenced papers,
whose citations are in the cover letter. The second phase is the actual
determination of degree of additional contribution by the reviewers. This
policy proposes that should a reviewer have questions, the AE is permitted
to reveal the citation to the 
reviewer (thus partially unblinding the paper), so that the reviewer can
fully check that the novelty requirement has been fulfilled.

Note however that this has the effect of devolving perhaps a good
percentage of reviews of such papers to SBR. However, the
author identity will still be omitted from the cover page of the submitted
paper, reminding the reviewer that the identity of the author should not
influence their review. The justification is that comprehensiveness of the
review is judged to be more important than blinding efficacy.

\subsubsection{On Revision}
The author guidelines includes a section on anonymity in revisions
(Appendix~\ref{app:rev}).

\noindent
{\sffamily Rationale}

Until the paper is accepted, author identities must remain masked.

\subsubsection{Acceptance}
As the AE is considering whether to accept the manuscript for publication,
they will consider the novelty requirements, which for double-blinded
reviews necessarily involve anonymous citations, which as noted above must
be fully documented in the cover letter.

If and when the manuscript is accepted for publication, the author will
prepare a non-anonymized version.
The AE will at that time make a final check to ensure that the \tods\
disclosure requirements are met.

\noindent
{\sffamily Rationale}

The specifics of meeting the disclosure requirements can be checked with the
final version of the paper. If there are any concerns, they can be
easily dealt with at that point.

\subsubsection{Transition}
It is important to educate the community as to the benefits and costs
of DBR. It is also important to publicize and promulgate policy changes
widely, {\em before} those changes go into effect. Finally, it is critical
to ensure that ManuscriptCentral is handling things correctly before
adopting such policy changes.

\noindent
{\sffamily Rationale}

The transition to DBR is a significant one for a journal, and so must be
done carefully. Such changes require time for the community to understand
and assimilate. The experience from the {\em SIGMOD} conference is that it
can take years for the community to come to embrace DBR.

\subsection{Cost-Benefit Analysis}

We revisit here the benefits and costs listed in Sections~\ref{sec:benefits}
and \ref{sec:costs}, respectively.

\begin{description}
\item [DBR avoids a bias towards top-ranked authors and institutions] The
experimental evidence is mixed concerning status bias. There is no clear
advantage to top-ranked authors and to authors from top-ranked institutions of
SBR, nor a clear disadvantage of DBR.

\item[DBR avoids a bias towards prolific authors] The contradictory results
of prior studies render it impossible to say anything definitive about the
impact of blinding on prolific authors.

\item[DBR encourages higher-quality reviews] It has not been shown
convincingly that either SBR or DBR can increase the quality of reviews of a
submitted manuscript.

\item[DBR is more fair to women authors] As all submissions are under DBR,
gender bias should be less of a concern.
\item[DBR is more fair to non-top-ranked authors and institutions] Similarly,
status bias should be less of a concern.
\item[There is a perception by some that DBR is more fair] Requiring DBR of
all submissions sends a strong message to the community that \tods\ is
making a sincere attempt to fully meet ACM's promise of an appropriate and
timely decision based on proper review by well-qualified and impartial
reviewers.
\end{description}

We now turn to an analysis of the costs of DBR of implementing the policy
just specified.

\begin{description}
\item[Quality of reviews is reduced] It has not been shown definitively that
DBR reduces the quality of reviews.

\item[Intentional bias is reduced] While there is substantial evidence of
referee bias in SBR, there is little evidence that referees were
intentionally biased in carrying out their duties as reviewers. Rather, the
bias that was observed seems to be implicit and unintended.

\item[It is difficult to mask submitted manuscripts]
With author masking, the source of the manuscript is available and the
blinding is done by the person most familiar with that manuscript, that is, the
author. Detailed guidelines and examples are provided in the author
instructions. It is anticipated that blinding a manuscript should take no
more than an hour of work.

We acknowledge that, in the presence of web search engines, submitted
manuscripts will be successfully masked less often than before such search
engines were available. The proposed approach does not require nor depend
on complete masking.

\item[DBR increases the probability of non-novel papers getting accepted]
In the reviewing procedure described here, reviewers are free to search for
related work. The responsibility for determining the contribution of the
work remains, whether the reviewing is single- or double-blind.

\item[It is difficult to fully mask systems papers]
The required anonymization steps are easy to apply to all papers. Full
masking is indeed difficult in some cases, and so is not required. All that
is required is to follow the simple six steps listed in the author instructions.

\item[DBR jeopardizes ``power projects'']
The submitted manuscript is of course free to mention the project upon which
it is based. What is not permitted is to state or imply that an author of
the submitted manuscript is one of the developers of that original
project.

As far as the reviewer knows, it is possible that another group has acquired the
software and has extended it, as exemplified in the author instructions.
Perhaps the author is a new graduate student in the research group,
or a student who has graduated and is at a different institution, or is a
friend of someone in the research group and was thus able to acquire the
system through that means.

All that is necessary is to introduce uncertainty as to the identity of the
author. 

Of course, the author will be clearly identified when the paper is
published, thereby extending and strengthening the ``brand'' of these power
projects.

\item[Forward-looking systems work, in particular, suffers from
anonymization]
It is recognized that some submissions reporting work on forward-looking
systems may not be truly anonymous.

Recall that the objective is to introduce uncertainty, not preclude all
possible ways of guessing the author.

In summary, it seems that the proposed process does not in any way put
systems papers, papers on power projects, or forward-looking papers at a
disadvantage. They are allowed to compete with other papers on an equal
footing; anonymization that obscures the scientific contributions of such
papers (or any other papers, for that matter) is not required.

\item[Anonymization is a distraction to students learning to write  
clearly]
The author guidelines give simple steps for converting a clearly
written paper into one for which the authorship is not revealed. These
six steps can be performed in the last hour before a manuscript is
submitted to \tods.

This requirement provides a useful pedagogical
opportunity to explain to students the vagaries of bias, the subtleties of
the peer review process, and finally, the need for the
{\em paper itself}, and not extraneous aspects such as the identity or
institution of the author(s), to convey the correctness of the arguments.

\item[Some journal policies preclude DBR]
The proposed policy shows that DBR can be executed while not jettisoning any
existing editorial policies for \tods.

\item[DBR precludes expanding a conference paper into a journal submission]
The proposed policy fully supports submissions expanded from published
conference papers.

\item[DBR makes it difficult for conflicts of interest to be detected]
The proposed process makes conflicts of interest explicit in the cover
letter. There is an initial cost to the author of preparing such a list. But
this cover letter helps to reduce the incremental load for the AEs, as
they would not have to avoid such conflicts only through their knowledge of
the network of relationships in the research community. (This is an
advantage regardless of the SBR-DBR distinction.)

\item[DBR prevents the dissemination of research results before submission]
The proposed policy dictates no constraints on such dissemination.
The only requirement is that anonymous citations be disclosed to the EiC and
AE in the cover letter (not to be shown to the reviewers).

\item[DBR gets papers rejected for failing to sufficiently cite the author's
own work]
\tods\ has specific novelty and disclosure
requirements that build on the ACM self-plagiarism
policy.\footnote{\url{http://www.acm.org/pubs/sim_submissions.html}} Not
following these guidelines could lead to rejection of the submitted
manuscript.

The anonymization guidelines in the author instructions show how the
author's own work may be cited. Doing so allows the author to simultaneously
satisfy the DBR, novelty, and disclosure requirements. The proposal mechanism
states when each of these requirements is
evaluated and ensures that adequate information is available at the
appropriate time.

\item[DBR increases the possibility of plagiarism] In the proposed
mechanism, reviewers are still responsible for detecting plagiarism. As the
community becomes more fluent with DBR, reviewers will not assume that they
know who all the authors are.

\item[DBR raises the possibility of self-plagiarism] The above analysis and
that in
Section~\ref{sec:assessing} argue that there are still effective processes
in place for detecting self-plagiarism. 

\item[DBR places extra burdens on reviewers] The six steps to blind a paper
shouldn't affect the readability of the paper.

Novelty is somewhat harder to check for a blinded submission. Should the
reviewer wish to check a reference that has been anonymized, they can ask
questions of the AE, who is permitted if needed to reveal the citation to
that reviewer.

\item[DBR decreases enthusiasm of reviewers]
There is little evidence from the literature that a policy of DBR reduces
incentives to do reviewing, reduces respect for reviewers, or makes it more
difficult for a journal to find reviewers. There may be an initial transient
effect when DBR is first adopted, but that should decrease as the community
becomes comfortable with such a policy.

\item[DBR reduces timeliness]
A previously-proposed protocol for checking novelty, requiring
several interactions between a reviewer and an AE, has now been simplified:
if needed, the reviewer just asks the AE for the full citation of an
anonymous citation.

\item[DBR places extra burdens on the EiC and AEs]
In the proposed process, the EiC has to configure MC so that it doesn't
reveal the authors to the reviewers in its dynamic web pages or email
templates. The EiC also must convey the new procedures to the community and
answer the inevitable questions. Finally, the EiC must examine each
submission with a quick scan to ensure that it has been properly blinded.

DBR also imposes additional work on AEs: they have to respond to questions
from reviewers regarding the material referenced by anonymous citations. On
the other hand, the list of conflicts provided by the authors makes
assigning reviewers slightly easier on AEs.

\item[DBR places extra burdens on the authors]
The authors are required to spend perhaps an hour following the steps in the
author instructions to anonymize their paper, as well as perhaps an hour
preparing the cover letter with the list of conflicts of interest, if that
list is not already available.

\item[SBR versus DBR should not be a local decision]
As the previous section noted, the prevalence of DBR has
increased dramatically over the last fifteen years, to the point where most
scientific journals now employ double-blind reviewing.
Closer to home, an informal vote at the SIGMOD 2006 business meeting in
Chicago indicated wide-spread support for DBR. A survey of database
researchers~\cite{Sirbu06} listed in the final recommendations ``To
maintain the perception of fairness (even though it might not have a large
impact on the acceptance rates), double-blind review could be
useful.''~[page 58].

Perhaps consistency is a benefit of DBR rather than a cost.

Looking at this another way, at the present time a database researcher has
absolutely no choice, as none of the journals currently employ mandatory
DBR. Having \tods\ adopt DBR provides choice to authors.
\end{description}

We have analyzed over twenty possible costs to DBR. Many of these costs
are completely eliminated through provisions of the proposal; examples
include jeopardizing power projects and forward-looking systems, journal
policies inconsistent with DBR, precluding expanding conference papers into
\tods\ submissions, and preventing dissemination of results before (or
after) submission. It seems that the primary costs are additional burdens on
the EiC, AEs, and authors. The burdens on the EiC and AEs are justified by
the desire of \tods\ to provide a reviewing process that is as fair as
possible. The burden on the authors, one or two hours per submission, is the
price for ensuring that gender bias, status bias, or other bias is less
likely to be a factor in the review of their submission.

We see that the administrative cost can be lowered through a combination of
author anonymization, system support (via MC), editor judgment in specific
cases (by the EiC for initial submission, and by the handling AE for
anonymous citations), and carefully-worded guidance that provides a specific
set of steps that can easily be executed by authors and reviewers.

\section{Discussion}
(While the literature survey and the cost-benefit analysis above attempt to be
objective, the present section presents the author's personal view.)

As the noted statistician Lynne Billard, who has written extensively on this
topic, has remarked, ``The issue of double-blind refereeing today is one
fraught with emotional overtones both rational and irrational, often
subconsciously culturally based, and so is difficult for many of us to
resolve equitably no matter how well intentioned''~\cite[page
320]{billard93}.

%TODO: comment from prominent, comment from Asian

\subsection{Why does this matter to \tods?}

The ACM Publications Board approved several years ago a broad rights and
responsibilities policy~\cite{ACMPubs,SnodgrassCACM}. Included was the
promise that ``The aim of the review process is to make an appropriate and
timely decision on whether a submission should be published. Such decisions
are based on proper review by well-qualified and impartial
reviewers. $\ldots$ Thus authors can expect ACM to use impartial reviewers
and to issue timely review and clear feedback.''

Each year the \tods' associate editors render editorial
decisions for about one hundred submitted manuscripts. The editors and
reviewers are for the most part diligent and careful in carrying out
the editorial process. That said, the scientific studies reviewed here and in the
companion paper suggest that for a small percentage of those submissions, the
wrong editorial decision was made.

It should be emphasized that {\em which} papers were erroneously rejected
(or accepted) is not known.  Referees never write ``this submission can't be
good: look where it came from.'' They don't even think that
consciously.

The impetus for these incorrect editorial decisions lies in an inherent bias in
the reviewing process that \tods\ uses~\cite{kardes05,wilson94}. This process provides irrelevant,
distracting information to reviewers that causes these reviewers, through no
intention or conscious effort, to see the paper in a more negative light.

No study has been done of the fairness of any ACM journal or
transaction. But very careful studies have been done of similar
journals. These studies have found systematic bias in unmasked
reviewing. The probability of such wrong decisions is evident only because
reviewer bias is so well-documented in the literature.

That this has occurred is not the fault of the \tods\ volunteer editors. It is
not the fault of the hundreds of reviewers, who are professional, caring,
objective scholars. It is the fault of a {\em process} that is patently
unfair, due to the various ways that bias psychologically affects judgment.

An especially productive way to avoid mental contamination of the judgment
process is to control one's exposure to biasing information, termed {\em
exposure control}~\cite{wilson94}. In terms of scholarly review, exposure
control is accomplished by blinding the reviewer to the identity and
affiliation of the author.

The literature also shows that ACM is an outlier in the scientific
community, in that it uses SBR for all of its journals. The majority of
scientific journals now use DBR. And
it is clear that the more a scientific discipline studies bias in general,
the more it utilizes DBR.

\subsection{Does bias exist in computer science?}
Bias in judgment is inevitable in complex tasks with many
dimensions~\cite{wilson94}. Journal reviewing is a prime example of such a
task. Major studies, described in Section~\ref{sec:fairness}, found evidence
of bias in economics and medical research. It is true that some submissions
to \tods\ are highly mathematical and thus might be less prone to reviewer
bias. However, the review of even highly mathematical papers requires subjective
judgment as to the contribution of the paper, a judgment in which reviewer
bias is certainly possible.

In the absence of a compelling countervailing mechanism, a reasonable
conclusion is that the \tods\ reviewing process, like that of the journals
for which those studies were done, is systematically unfair, despite
the promises of the afore-mentioned ACM policy.

\subsection{What is the appropriate decision procedure?}\label{sec:decisionprocedure}

The ACM rights and responsibilities policy puts the onus on ACM and on the
EiC to ensure that the reviewing process is fair.

Thus the issue is not, how to make the reviewing fair when it is easy
to administer. \tods\ goes to great lengths, at great cost, to ensure
fair reviewing. Reviewer identity is held in confidence. The review form
lists several aspects to consider, to guide the reviewer in considering
appropriate aspects. Reviews are much more than just recommendations
(accept/reject): reviewers are expected to document in detail their
rationale for their recommendation.  Several reviewers are used for
each paper, to not be overly reliant on just one or two views. These
elaborate procedures require months per submission. An editorial
decision could be made much more easily, in just a few days or weeks,
if fairness wasn't viewed as critical.

The appropriate decision procedure is to first decide, is
SBR or DBR more fair? (This is assuming that DBR can be done at all, but
many journals have shown by example that DBR is possible with journals.) And
only then, if DBR is judged to be more fair, to determine how the process
can be made as smooth as possible and with as little cost as possible, while
retaining fairness and balancing remaining costs based on stated principles.

Some of the costs can be reduced unilaterally, as shown in
Section~\ref{sec:proposal}. For example, journal policies that preclude DBR
can be adjusted to accommodate DBR and conflicts of interest can be
identified even with DBR.

Some of the costs are discipline-specific. For disciplines that do not
involve ``power projects,'' such as the more mathematically-oriented
disciplines, a provision for ambiguating statements on well-known or unique
systems may not be relevant. For disciplines in which there is not a
tradition of extending conference papers into journal submissions, a
procedure for revealing anonymous citations may not be needed. These are
``knobs'' that a policy designer can adjust to fine-tune the policy.

But many of the costs, if reduced, increase other costs, due to complex
interactions. The approach taken in Section~\ref{sec:proposal} was to apply
a principled tradeoff between such costs and blinding efficacy. These are
the other ``knobs'' that are available to policy designers.

Three principles are used here. The first is that authors should not be
required to go to great lengths to blind their submissions. The second is
that comprehensiveness of the review trumps blinding efficacy. The final
principle used here is that AEs retain flexibility and authority in managing
the reviewing process.  Hence, the procedure dictates six quick steps for
blinding a submission, when a more aggressive obfuscation of the paper would
yield higher blinding efficacy. And the procedure explicitly allows the AE to
reveal author identity information, such as a full citation that had been
previously anonymized, if needed by a reviewer to evaluate the submission.

Applying these principles encouraged a systematic design of the tradeoff
between costs and blinding efficacy.

\subsection{Decision}

In considering this decision, it is my responsibility as Editor-in-Chief to
understand the scientific data and to gather input from all
constituencies. I have read most of the papers concerning blind review and
have talked to perhaps one hundred people.

The data shows that status bias is against those in the gray area:
neither at the top (prominent, prolific, or from top schools: the
{\em scholarly elite}) nor at the bottom (those who can't write a good
paper). It is those in the gray area who are hurt by the current
system.

I involved the \tods\ editorial board because it is the AEs who must execute
the submission handling process.  I worked with the editorial board for two
months to make the process as smooth as possible. The AEs are the experts on
the process and have provided many excellent suggestions that have resulted
in the clean, workable design presented in Section~\ref{sec:proposal}.

Concerning the scholarly elite, I consider the \tods\ editorial board to be a
representative sample. The consensus among the \tods\ editorial board is
that SBR is less fair than DBR. When asked for their recommendation,
roughly a third recommended SBR, roughly a third recommended
DBR, and roughly a third were comfortable with either SBR or DBR.

But the primary constituency consists of qualified potential authors,
those who write reasonable papers but who suffer from status bias from
a flawed process. Note that none of these people is on the editorial board.

An informal vote taken at the SIGMOD business meeting in Chicago in June
2006, in a room
filled with hundreds of qualified potential authors, was decisive: over 90\%
were in favor of DBR. That vote is mirrored in the individual
discussions I have had. Most are astounded that there is even a
question.

Given this information, my decision was a natural one. The many
qualified potential authors, for whom SBR may be less fair, overwhelmingly
favor DBR. The scholarly elite, who are small in number and who might in
some cases unfairly benefit from the current process, are mixed. And it has
been shown to be possible to design an editorial process that
minimizes the costs while still emphasizing the primacy of the merits of the
submission.

My conclusion is simple: the \tods\ reviewing process should be fair to
everyone.

\section{TODS Policy}

The following is now \tods\ policy.
\begin{itemize}
\item \tods\ reaffirms the general ACM policy that ``the quality of a refereed
  publication rests primarily on the impartial judgment of their volunteer
  reviewers.''

\item \tods\ will continue to strive to ensure fairness in reviewing, even if
  that involves more work for the \tods\ editorial board.

\item Scientific studies have demonstrated opportunities for bias inherent in
  single-blind reviewing.

\item It is \tods\ policy that every submission should be judged on its own
merits. The identity and affiliation of the authors should not influence,
  either positively or negatively, the evaluation of submissions to \tods.

\item In consideration of the above, \tods\ will utilize double-blind reviewing.

\item \tods\ will continue to strive to make the submission process for
  authors as simple as possible.

\item \tods\ will continue to strive to effect a comprehensive review of each
submission.
\end{itemize}

This policy is not dependent on absolute or even relative blinding
efficacy. The central and unambiguous message is that every submission
should be judged solely on its own merits. This message applies even when
reviewers know exactly who the authors are.
The other important message is that \tods\ so values fairness that it is
willing to undertake additional effort by AEs to make the process
more fair.

Given that the community has spoken with such a clear
voice and that there are so few instances of journals in which mandatory DBR
has been subsequently replaced with SBR, I believe that
this conclusion will be embraced by the vast majority of the
community, as authors and reviewers and AEs gain experience with the
reviewing process.  To aid in this transition, I have prepared a brief
FAQ~\cite{FAQ}.

\section*{Acknowledgments}
The instructions to authors and instructions to reviewers given in the
appendix are based on instructions given to
{\em SIGMOD'06} authors. These instructions have benefited from editing by
the program chairs of {\em SIGMOD'01} through {\em SIGMOD'06}: Timos Sellis,
Michael Franklin, Alon Halevy, Gerhard Weikum, Jennifer
Widom, and Surajit Chaudhuri.

The author thanks the many \tods\ associate editors and others who have
corresponded privately with their ideas and concerns about SBR and DBR,
including but not limited to Kurt Akeley, Michel Beaudouin-Lafon, Alan
Berenbaum, Ronald Boisvert, Judy Brown, Merrie Brucks, Boots Cassel, Surajit
Chaudhuri, Jan Chomicki, Douglas D.~Dankel II, Peter Denning, David DeWitt,
Susan Dumais, Curtis Dyreson, Kemal Ebcioglu, Carla Ellis, Mary Fernandez,
Michael Franklin, Richard P.~Gabriel, Luis Gravano, Ralf G\"{u}ting, Alon
Halevy, Joe Hellerstein, Mary Jane Irwin, Christian S.~Jensen, Susan Jones,
Panagiotis Karras, Won Kim, Phokion Kolaitis, Hank Korth, Donald Kossmann,
David Leake, Steven Levitan, Alan Mackworth, David Maier, Kathryn McKinley,
Elli Mylonas, Fred Niederman, Beng Chin Ooi, Z.~Meral \"{O}zsoyo\v{g}lu,
Tamer \"{O}zsu, Craig Partridge, Ray Perrault, Michael Pristley, Raghu
Ramakrishnan, Stephanie Rosenbaum, Kenneth Ross, Holly Rushmeier, Betty
Salzberg, Pierangela Samarati, Sunita Sarawagi, Stuart Selber, David
E.~Siegel, Mary Lou Soffa, Dan Suciu, S.~Sudarshan, Geoff Sutcliffe, Marilyn
Tremaine, Jan Van den Bussche, Mary Vernon, Jennifer Widom, Nina Wishbow,
and Xindong Wu.  Their comments have greatly improved this analysis.

\addcontentsline{toc}{section}{References}
%\bibliographystyle{alpha}
\newcommand{\etalchar}[1]{$^{#1}$}
\begin{thebibliography}{}
%\begin{thebibliography}{BADW82}
%\begin{thebibliography}{[99]}

\bibitem[ACM 2001]{ACMPubs}
{\sc ACM Publications Board}, Rights and Responsibilities in ACM Publishing,
approved June 27, 2001. {\tt http://www.acm.org/pubs/rights.html}, viewed
December 5, 2006.

\bibitem[Altman et al.~1991]{altman91}
%\bibitem{altman91}
{\sc N.~Altman, D.~Banks, P.~Chen, D.~Duffy, J.~Hardwick, C.~L\'{e}ger, A.~Owen, and
T.~Stukel}, Meeting the Needs of New Statistical Researchers, {\em
  Statistical Science} 6(2):163--174, May 1991.

\bibitem[Altman et al.~1992]{altman92}
%\bibitem{altman92}
{\sc N.~Altman, J.~F.~Angers, D.~Banks, D.~Duffy, J.~Hardwick, C.~L\'{e}ger,
M.~Martin, D.~Nolan (Chair), A.~Owen, D.~Politis, K.~Roeder, T.~N.~Stukle
and Z.~Ying}, Rejoinder, {\em Statistical Science} 7(2):265--266, May 1992.

\bibitem[Billard 1993]{billard93}
%\bibitem{billard93}
{\sc L.~Billard}, Comment, {\em Statistical Science} 8(3):320--322, August
1993.

\bibitem[Calfee \& Valencia 2006]{APAGuide}
%\bibitem{APAGuide}
{\sc R.~C.~Calfee and R.~R.~Valencia}, APA Guide to Preparing Manuscripts for
Journal Publication, \url{http://www.apa.org/journals/authors/guide.html},
viewed June 21, 2006.

\bibitem[Blank 1991]{blank91}
%\bibitem{blank91}
{\sc R.~M.~Blank}, The Effects of Double-Blind versus Single-Blind Reviewing:
Experimental Evidence from The American Economic Review, {\em American
  Economic Review} 81(5):1041--1067, December 1991.

\bibitem[Borts 1974]{borts74}
%\bibitem{borts74}
{\sc G.~Borts}, Report of the Managing Editor, {\em American Economic Review}
64:476--482, May 1974.

\bibitem[Ceci \& Peters 1984]{ceci84}
%\bibitem{ceci84}
{\sc S.~J.~Ceci and D.~P.~Peters}, How blind is blind review? {\em American
  Psychologist} 39:1491--1494, 1984.

\bibitem[Cox et al.~1993]{cox93}
%\bibitem{cox93}
{\sc D.~Cox, L.~Gleser, M.~Perlman, N.~Reid, and K.~Roeder}, Report of the Ad Hoc
  Committee of Double-Blind Refereeing, {\em Statistical
  Science} 8(3):310--317, August 1993.

\bibitem[Dalton 1995]{dalton95}
%\bibitem{dalton95}
{\sc M.~Dalton}, Refereeing of Scholarly Works for Primary Publishing, in {\em
  Annual Review of Information Science and Technology (ARIST)}, Volume 30,
{\sc M.~E.~Williams} (ed.), American Society of Information Science, 1995.

\bibitem[Garfunkel et al.~ 1994]{garfunkel94}
%\bibitem{garfunkel94}
{\sc J.~M.~Garfunkel, M.~H.~Ulshen, H.~J.~Hamrick, and E.~E.~Lawon}, Effect of
Institutional Prestige on Reviewers' Recommendations and Editorial
Decisions, {\em Journal of the American Medical Association}
272(2):137--138, July 13, 1994.

\bibitem[Genest 1993]{genest93}
{\sc C.~Genest}, Comment, {\em Statistical Science} 8(3):323--327, August,
1993.

\bibitem[Gordon 1980]{gordon80}
{\sc M.~D.~Gordon}, The role of referees in scientific communication, in {\em
The Psychology of Written Communication: Selected Readings}, {\sc J.~Hartley}
(ed.), London, England: Kogan Page, pp.~263--275, 1980.

\bibitem[Greenwald \& Banaji 1995]{greenwald95}
{\sc A.~G.~Greewald and M.~R.~Banaji}, Implicit social cognition: Attitudes,
self-esteem, and stereotypes, {\em Psychological Review}, 102(1):4--27,
January 1995.

\bibitem[Horrobin 1982]{horrobin82}
%\bibitem{horrobin82}
{\sc D.~F.~Horrobin}, Peer review: A philosophically faulty concept which is
proving disastrous to science, {\em Behavioral and Brain
Sciences} 5(2):217--218, 1982.

\bibitem[Kardes et al.~2005]{kardes05}
{\sc F.~R.~Kardes, A.~V.~Muthukrishnan, and V.~Pashkevich}, On the
conditions under which experience and motivation accentuate bias in
intuitive judgment, in {\em The Routines of Decision Making}, {\sc
  T.~Betsch and S.~Haberstroh} (ed.), Lawrence Erlbaum Associates,
pp.~139--156, 2005.

\bibitem[Link~1998]{link98}
%\bibitem{link98}
{\sc A.~M.~Link}, US and Non-US Submission: An Analysis of Review Bias, {\em
  Journal of the American Medical Association} 280(3):246--247, July 15, 1998.

\bibitem[Mahoney et al.~1978]{mahoney78}
{\sc M.~J.~Mahoney, A.~E.~Kazdin, and M.~Kenigsberg}, Getting Published, {\em
Cognitive Therapy and Research} 2(1):69--70, 1978.

\bibitem[Montesi 2006]{Montesi06}
{\sc M.~Montesi}, From Conference to Journal Publication. How Conference Papers
in Software Engineering are Extended for Publication in Journals, in
preparation, 2006.

\bibitem[NSF 2004]{NSF}
{\sc National Science Foundation}, Conflict-of-Interests and Confidentiality
Statement for NSF Panelists, NSF Form 1230P, February 2004.

\bibitem[Rosenblatt \& Kirk 1980]{rosenblatt80}
{\sc A.~Rosenblatt and S.~A.~Kirk}, Recognition of Authors in Blind Review of
Manuscripts, {\em Journal of Social Service Research} 3(4):383--394,
Summer 1980.

\bibitem[Shiv et al.~2005]{shiv05}
{\sc B.~Shiv, Z.~Carmon, and D.~Ariely}, Placebo Effects of Marketing
Actions: Consumers May Get What They Pay For, {\em Journal of Marketing
  Research} 42(4):383--393, November 2005.

\bibitem[Sirbu and Vasilescu 2006]{Sirbu06}
{\sc A.~Sirbu and A.-T.~Vasilescu}, The Health of Database Research, Project
Report, Worcester Polytechnic Institute, May 2006.

\bibitem[Snodgrass 2002]{SnodgrassCACM}
{\sc R.~T.~Snodgrass}, Rights and Responsibilities in ACM Publishing,
{\em CACM} 45(2): 97--101, February 2002.

\bibitem[Snodgrass 2005]{AEmanual}
{\sc R.~T.~Snodgrass}, {\em {\em ACM TODS} Associate Editor Manual},
September 2005.

\bibitem[Snodgrass 2006]{DBR}
{\sc R.~T.~Snodgrass}, ``Single- Versus Double-Blind Reviewing: An Analysis of
the Literature,'' {\em ACM SIGMOD Record}, 35(3):8--21, September 2006.

\bibitem[Snodgrass 2007]{FAQ}
{\sc R.~T.~Snodgrass}, ``Frequently-Asked Querstions About Double-Blind Reviewing,'' {\em ACM SIGMOD Record}, 36(1), March 2007.

\bibitem[van Rooyen 1999]{vanRooyen99}
%\bibitem{vanRooyen99}
{\sc S.~van Rooyen, F.~Godlee, S.~Evans, R.~Smith, and N.~Black}, Effect of
Blinding and Unmasking on the Quality of Peer Review: A Randomized Trial,
{\em Journal of General Internal Medicine} 14(10):622--624, October 1999.

\bibitem[Werner 2006]{werner06}
{\sc S.~Werner}, personal communication, December 4, 2006.

\bibitem[Wilson \& Brekke 1994]{wilson94}
{\sc T.~D.~Wilson and N.~Brekke}, Mental contamination and mental
correction: Unwanted influences on judgments and evaluations, {\em
  Psychological Bulletin} 116(1):117--142, July 1994.

\bibitem[Zuckerman \& Merton 1971]{zuckerman71}
%\bibitem{zuckerman71}
{\sc H.~Zuckerman and R.~K.~Merton}, Patterns of Evaluation in Science:
Institutionalization, Structure and Functions of the Referee System,
{\em Minerva} 9(1):66--100, January 1971.

\end{thebibliography}

\appendix
\section{Instructions to Authors}\label{app:authors}

\tods\ strives to ensure fairness in reviewing. It is
\tods\ policy that every submission should be judged solely on its own
merits. The identity and affiliation of the authors should not influence,
either positively or negatively, the evaluation of submissions to \tods.
For this reason, all submissions to \tods\ will undergo {\em double-blind
review}, in which authors and reviewers are unaware of each other's
identities.

\subsection{Ensuring Anonymity}
To ensure anonymity of authorship, authors must blind their manuscript by
performing the following simple alterations.
\begin{itemize}
\item  Authors' names and affiliations must not appear on the title page or
elsewhere in the paper. 
\item Funding sources(s) must not be acknowledged on the title page or
elsewhere in the paper. 
\item All personal acknowledgments should be omitted. Research group members or other
colleagues or collaborators must not be acknowledged anywhere in the
paper. There should also be no acknowledgment section in the paper. 
\item Source file naming must also be done with care. It is strongly
suggested that the submitted file be named as {\tt todssubmission.pdf}. Also,
if your 
name is Jane Smith and you submit a PDF file generated from a .dvi file
called {\tt Jane-Smith.dvi}, one can infer your authorship by looking into the
PDF file. Similarly, you should remove name and affiliation information from
the ``properties'' of your document.  Some forget this, since such properties
are filled in silently when a document is created or copied.
For example, Microsoft Word fills in your Name and Company that it
maintains from the installation process.  This information
is sometimes transferred even when you covnert to another
form, such as PDF.
\end{itemize}

Despite the anonymity requirements, you should still include relevant prior
published work of your own in the references---omitting them could
potentially reveal your identity by negation. {\em Prior published work} is
defined as any research paper that has been published and made available
prior to submission to \tods\ (a)~in the online or printed proceedings of a
refereed conference or refereed workshop, (b)~as a longer poster papers (4
pages or more) in such a conference or workshop, or (c)~in an online or
printed issue of a journal.  However, you should not cite a published
demonstration paper of your own even if it is in prior published work as
defined above.

You must use care in referring to your prior published work. For
example, if you are Jane Smith, the following text reveals the authorship
of the submitted paper:

\pagebreak

\begin{quotation}
In our previous work [Doe 1997, Smith 1998], we presented two algorithms for .... We build on that work by ...

{\sffamily Bibliography}

{\sc John Doe and Jane Smith}, A Simple Algorithm for ..., {\em Proceedings of ACM
SIGMOD International Conference on Management of Data}, pp.~1--10, 1997. 

{\sc Jane Smith}, A More Complicated Algorithm for.., {\em Proceedings of ACM
SIGMOD International Conference on Management of Data}, pp. 34--44, 1998.
\end{quotation}

\noindent
The solution is to reference your prior published work in the third person
(just as you would any other piece of work that is related to the submitted
paper). This allows you to set the context for the submitted paper, while at
the same time preserving anonymity.

\begin{quotation}
In previous work [Doe 1997, Smith 1998], algorithms were presented for ... We
build on that work by ...

{\sffamily Bibliography}

{\sc John Doe and Jane Smith}, A Simple Algorithm for ..., {\em Proceedings of ACM
SIGMOD  International Conference on Management of Data}, pp.~1--10, 1997.

{\sc Jane Smith}, A More Complicated Algorithm for.., {\em Proceedings of ACM
SIGMOD  International Conference on Management of Data}, pp. 34--44, 1998.
\end{quotation}

\subsection{Anonymous Citations}\label{app:anoncit}
Referring to any work (including but not limited to your own work) that has
been submitted elsewhere for review (and hence is as yet unpublished), or
that has been accepted for publication at a refereed conference, refereed
workshop, or a journal for which proceedings (printed or online) will be
made available after the submission of your \tods\ manuscript, requires a
different protocol to ensure consistency with double-blind reviewing. In the
body of your \tods\ submission, you may refer to such work in the third
person as follows.

\begin{quotation}
The
authors have also developed closely related techniques for query
optimization [Anonymous], but...
\end{quotation}

\noindent
In the above example, the reference in the bibliography would then read:

\begin{quotation}
{\sc Anonymous}, details omitted due to double-blind reviewing.
\end{quotation}

Note that you should mention neither the authors nor the title or venue of
publication while describing {\em anonymous citations} like that above. Be
sure to place all anonymous citations after the list of your regular
citations. Of course, this does not mean you are not responsible to provide
the details of such anonymous citations. You should disclose (for use by the
Editors only) full details of each such anonymous citation in your cover
letter, which will not be made available to the reviewer. Furthermore, you
may be asked to submit a copy of one of these papers corresponding to such
citations. An editor will contact you during the review period if this
becomes necessary.

Technical reports (or URLs for down-loadable versions) of your own work
should not be referenced. Self-references should also be limited to only
papers that are relevant and essential for the reviewing of the
submitted paper.

\subsection{Conflicts of Interest}
As conflicts of interest between the author and reviewer are harder to
detect in double-blind review, authors are required to submit a list of all such conflicts in their
cover letter. To reduce the effort of preparing such a list, \tods\ adopts the
definition of conflict of interest utilized by the US National Science
Foundation.

Specifically, you should list all people who are potential reviewers and who
have a relationship with one of the authors involving the following.
\begin{itemize}
\item Known family relationship as spouse, child, sibling, or parent.
\item Business or professional partnership.
\item Past or present association as thesis advisor or thesis student.
\item Collaboration on a project or on a book, article, report, or paper within the last 48 months.
\item Co-editing of a journal, compendium, or conference proceedings within the last 24 months.
\item Other relationship, such as close personal friendship, that you think
might tend to affect your judgment or be seen as doing so by a reasonable
person familiar with the relationship.
\end{itemize}

\noindent
Please list these people under the appropriate categories.

\subsection{Well-Known or Unique Systems}\label{app:sys}
Papers on well-known or unique systems may be more difficult to blind.
All that is required is that blatant statements that could identify one of
the authors be ambiguated.

A common approach is to state that the developers of the system (which might
well have been the authors of the submitted paper) made the system available
to the authors of the submitted paper (which is strictly a true
statement). For software, the developers might have shared the source code
with the authors; for a hardware system, the developers might have made the
system available for use by the authors. It might even be useful to
explicitly mention that the developers provided guidance on the
internals of the system.

There is no expectation that manuscripts be extensively rewritten to attempt
to mask the identity of the authors.

\subsection{Anonymity in Revisions}\label{app:rev}
All comments to the reviewers and all manuscript revisions except for a
possible final revision after the paper has been accepted (should the paper
be accepted) should also follow the guidelines above.

\subsection{Summary}
It is the responsibility of authors to do their very best to follow the
specific steps listed here to preserve anonymity. Papers that do not follow
the guidelines here are subject to immediate rejection.

The steps can be summarized quickly.
\begin{enumerate}
\item Anonymize the title page.
\item Remove mention of funding sources and personal acknowledgments.
\item Anonymize references found in running prose that cite your papers.
\item Anonymize citations of submitted work in the bibliography.
\item Ambiguate statements on well-known or unique systems that identify an author.
\item Name your files with care and ensure document properties are also anonymized.
\end{enumerate}

\noindent
Authors need only take these six steps to adequately blind their papers.

Common sense can go a long way toward preserving
anonymity without diminishing the quality or impact of a paper. The goal is
to preserve anonymity to a reasonable degree while still allowing the reader
to fully grasp the context (related past work, including your own) of the
submitted paper and while making it relatively easy for the author to follow
the required steps. If you desire specific guidance, please contact the
Editor-in-Chief.

\section{Instructions to Reviewers}\label{app:reviewers}

It is \tods\ policy to use {\em double-blind reviewing} for all papers, in
which authors and reviewers are unaware of each other's identities.
It is \tods\ policy that every submission should be judged solely
on its own merits. The identity and affiliation of the authors should not
influence, either positively or negatively, the evaluation of submissions to
\tods. In particular, you are asked not to go
to unusual lengths to try to discover the identity of the author. On the
other hand, you should still be diligent in determining the contribution of
the submission over previously-published work by the author or by others.

Authors have been given guidelines on how to approach double-blind
reviewing. Please review the submission page to familiarize yourself with
these guidelines. If authors have not made reasonable effort to conform to
double-blind reviewing, the paper could be rejected.  ``Reasonable effort''
means the author followed the guidelines in the submission page. Your
ability to guess/discover an author's identity based on your past
experience, extensive web search, non-obvious ``clues'' in the submitted
paper, and other exploratory means are not necessarily grounds for
rejection.

Should you need access to the material referenced by anonymous citations
(those stating ``details omitted due to double-blind reviewing''), say to
judge the novelty requirement, simply notify the Associate Editor handling
this paper and give your reason and the full citation will be revealed to
you. If you have any other questions about double-blind reviewing, please
also contact the Associate Editor.

\end{document}
